\documentclass[11pt,letterpaper]{article}

%%% LANGUAGE: Spanish
\usepackage[spanish,english]{babel}
\addto\shorthandsspanish{\spanishdeactivate{~<>}}

%%% ENCODING
\usepackage[utf8]{inputenc}
\usepackage{microtype}
\usepackage[T1]{fontenc}
\usepackage{lmodern}
\usepackage{amsmath}
\usepackage{textcomp}

\usepackage{graphicx}
\usepackage{epstopdf}
\usepackage{epsfig}

%%
%%%%GENERAL INFROMATION%%%%
%%

\newcommand{\theCampus}{Campus Vitacura}
\newcommand{\theAssignment}{SIG}
\newcommand{\therevision}{DV}

%%% PACKAGES
\usepackage{booktabs} 	% for much better looking tables
\usepackage{array} 		% for better arrays (eg matrices) in maths
\usepackage{paralist} 	% very flexible & customizable lists (eg. enumerate/itemize, etc.)
\usepackage{verbatim} 	% adds environment for commenting out blocks of text & for better verbatim
\usepackage{subfigure} % make it possible to include more than one captioned figure/table in a single float
\usepackage{xcolor}
\usepackage{listings}  % bloques de código SQL

% FECHA
\usepackage{datetime}
\newdateformat{monthyeardate}{%
  \monthname[\THEMONTH], \THEYEAR}

%%% PAGE DIMENSIONS
\usepackage[margin=25mm,inner=30mm]{geometry} % to change the page dimensions


%%% HEADERS & FOOTERS   %%%
\usepackage{fancyhdr} 	% This should be set AFTER setting up the page geometry
\pagestyle{fancy} 			% options: empty , plain , fancy
\renewcommand{\headrulewidth}{0.5pt} % customise the layout...

\renewcommand{\footrulewidth}{0.5pt} % customise the layout...
\lhead{\includegraphics[scale=.15]{Logos/header_utfsm.png}}\cfoot{}
\cfoot{\textbf{\thepage}}
\rfoot{\therevision}
\lfoot{\theAssignment}\chead{}
\rhead{}


\usepackage[colorlinks=true,linkcolor=blue,urlcolor=blue]{hyperref}

%%% SECTION TITLE APPEARANCE
\usepackage{sectsty}
%\allsectionsfont{\sffamily\mdseries\upshape} % (See the fntguide.pdf for font help)
\usepackage{float}
\usepackage{titlesec}
\titlelabel{\thetitle.\quad}

\usepackage[numbers]{natbib}

\usepackage{lipsum}
\usepackage[section]{placeins}

%%% ESTILO DE CÓDIGO SQL
\lstdefinestyle{sqlstyle}{
  language=SQL,
  basicstyle=\ttfamily\small,
  keywordstyle=\bfseries\color{blue!60!black},
  stringstyle=\color{orange!70!black},
  commentstyle=\color{gray},
  numbers=none,
  columns=flexible,
  keepspaces=true,
  upquote=true,
  showstringspaces=false,
  breaklines=true,
  frame=single,
  rulecolor=\color{gray!60},
  backgroundcolor=\color{gray!8},
  framesep=4pt,
  xleftmargin=2pt,
  xrightmargin=2pt,
  aboveskip=6pt,
  belowskip=6pt
}
\lstset{style=sqlstyle}

% Figura de captura de MySQL Workbench (ancho común para todas)
\newcommand{\captura}[3]{%
  \begin{figure}[H]
    \centering
    \includegraphics[width=0.88\linewidth]{Lab4/Figuras/#1}
    \caption{#2}
    \label{#3}
  \end{figure}
}


%------------------------------------------------------------------------------
%%%% TITLE PAGE
%------------------------------------------------------------------------------
\begin{document}
\selectlanguage{spanish}

\begin{titlepage}
\begin{figure}[H]
\includegraphics[width = 0.4\linewidth]{Logos/logo_di.png}
\hspace{144pt}
\vspace{-70pt}
\end{figure}
\vspace{60mm}
\begin{center}
\hrule
\vspace{2mm}
{\Huge \textbf{\sc Informe de Laboratorio 4:}} \\
\vspace{2mm}
{\Large \sc{\textit{Entregable SQL (1/2): base de datos de Austral Editores}}}
\vspace{2mm}
\hrule
\vspace{1cm}


%%%%AUTHORS%%%%
\centering
    {\textbf{Diego Victoriano Rodríguez} \\
    {\small Ingeniería Civil Industrial, 21.540.820-5, \emph{diego.victoriano@usm.cl}}\\
    }

    \vspace{0.5cm}
    {ICN292 - Sistemas de Información para la Gestión.} \\
    {\textbf{Paralelo:} 101}

    \vspace{0.5cm}
    \textbf{Profesores:} José Miguel González P. \textbar\ José Luis Sáez T. \textbar\ Patricio Cifuentes\\[0.3cm]
    \textbf{Ayudantes:} Valeria Bobadilla \textbar  Catalina Bodaleo \textbar\ Luis Contreras \textbar  Benjamín Jara

    \vspace{0.8cm}
    \textbf{Repositorio GitHub:}\\[0.1cm]
    {\small \url{https://github.com/diegovictorianor/ICN292-Lab4-Victoriano-Diego}}

    \vfill
    \textsc{\monthyeardate\today}\\[0.5cm]
    \textsc{\theCampus}
\end{center}


\clearpage

\parindent=0mm


\end{titlepage}

%--------------------------------------------------
%%% CONSULTAS %%%%
%--------------------------------------------------

%%% Consulta 1
\section*{Consulta 1: autores cuyo país no es Chile}

\begin{lstlisting}
SELECT nombre, genero
FROM autores
WHERE pais <> 'Chile' OR pais IS NULL
ORDER BY nombre;
\end{lstlisting}

La comparación \lstinline|pais <> 'Chile'| devuelve verdadero para Inglaterra, Argentina y Estados Unidos, falso para Chile y NULL (desconocido) cuando el país está vacío, como en Kenji Arakawa; como \texttt{WHERE} solo conserva las filas verdaderas, ese autor se perdería sin la condición adicional \texttt{OR pais IS NULL}.

\captura{consulta_1.png}{Resultado de la consulta 1 en MySQL Workbench.}{fig:c1}

\clearpage
%%% Consulta 2
\section*{Consulta 2: libros de 190 a 220 páginas}

\begin{lstlisting}
SELECT titulo, paginas
FROM libros
WHERE paginas BETWEEN 190 AND 220
ORDER BY paginas ASC;
\end{lstlisting}

\captura{consulta_2.png}{Resultado de la consulta 2 en MySQL Workbench.}{fig:c2}

\clearpage
%%% Consulta 3
\section*{Consulta 3: resumen del catálogo en una fila}

\begin{lstlisting}
SELECT COUNT(*)                 AS total_libros,
       SUM(ejemplares_vendidos) AS total_ejemplares,
       AVG(paginas)             AS promedio_paginas,
       MIN(anio)                AS anio_mas_antiguo,
       MAX(ejemplares_vendidos) AS max_ejemplares
FROM libros;
\end{lstlisting}

\captura{consulta_3.png}{Resultado de la consulta 3 en MySQL Workbench.}{fig:c3}

\clearpage
%%% Consulta 4
\section*{Consulta 4: géneros distintos en orden alfabético}

\begin{lstlisting}
SELECT DISTINCT genero
FROM autores
ORDER BY genero;
\end{lstlisting}

\captura{consulta_4.png}{Resultado de la consulta 4 en MySQL Workbench.}{fig:c4}

\clearpage
%%% Consulta 5
\section*{Consulta 5: autores con más de un libro publicado en 2024}

\begin{lstlisting}
SELECT a.nombre, COUNT(*) AS libros_2024
FROM autores a
JOIN libros l ON a.id_autor = l.id_autor
WHERE l.anio = 2024
GROUP BY a.id_autor, a.nombre
HAVING COUNT(*) > 1
ORDER BY a.nombre;
\end{lstlisting}

\captura{consulta_5.png}{Resultado de la consulta 5 en MySQL Workbench.}{fig:c5}

\clearpage
%%% Consulta 6
\section*{Consulta 6: libros de 2024 con más de 2\,000 ejemplares}

\begin{lstlisting}
SELECT titulo, ejemplares_vendidos
FROM libros
WHERE anio = 2024 AND ejemplares_vendidos > 2000
ORDER BY ejemplares_vendidos DESC
LIMIT 3;
\end{lstlisting}

\captura{consulta_6.png}{Resultado de la consulta 6 en MySQL Workbench.}{fig:c6}

\clearpage
%%% Consulta 7
\section*{Consulta 7: insertar y eliminar el libro \emph{Papel Mojado}}

\begin{lstlisting}
INSERT INTO libros (id_autor, titulo, paginas, ejemplares_vendidos, anio)
VALUES (3, 'Papel Mojado', 260, 50, 2024);

SELECT * FROM libros WHERE titulo = 'Papel Mojado';

DELETE FROM libros WHERE titulo = 'Papel Mojado';

SELECT * FROM libros WHERE titulo = 'Papel Mojado';
\end{lstlisting}

\captura{consulta_7.1.png}{Resultado de la consulta 7 (\texttt{INSERT} y \texttt{SELECT} previo al \texttt{DELETE}) en MySQL Workbench.}{fig:c7a}

\captura{consulta_7.2.png}{Resultado de la consulta 7 (\texttt{DELETE} y \texttt{SELECT} posterior) en MySQL Workbench.}{fig:c7b}

\clearpage
%%% Consulta 8
\section*{Consulta 8: sumar 100 ejemplares a \emph{Cielo de Enero}}

\begin{lstlisting}
SELECT titulo, ejemplares_vendidos
FROM libros WHERE titulo = 'Cielo de Enero';

UPDATE libros
SET ejemplares_vendidos = ejemplares_vendidos + 100
WHERE titulo = 'Cielo de Enero';

SELECT titulo, ejemplares_vendidos
FROM libros WHERE titulo = 'Cielo de Enero';
\end{lstlisting}

\captura{consulta_8.1.png}{Resultado de la consulta 8 (\texttt{SELECT} previo) en MySQL Workbench.}{fig:c8a}

\captura{consulta_8.2.png}{Resultado de la consulta 8 (\texttt{UPDATE} y \texttt{SELECT} posterior) en MySQL Workbench.}{fig:c8b}

\clearpage
%%% Consulta 9
\section*{Consulta 9: libros, colección y fecha de publicación de la colección}
\begin{lstlisting}
SELECT l.titulo AS titulo_libro,
       c.titulo AS titulo_coleccion,
       c.fecha_publicacion
FROM libros l
JOIN colecciones c ON l.id_coleccion = c.id_coleccion
ORDER BY c.fecha_publicacion, l.titulo;
\end{lstlisting}

\captura{consulta_9.png}{Resultado de la consulta 9 en MySQL Workbench.}{fig:c9}

\clearpage
%%% Consulta 10
\section*{Consulta 10: pagos realizados a los autores}

\subsection*{Sentencia 1: pagos hechos a los autores}

\begin{lstlisting}
SELECT a.nombre, p.fecha_pago, p.monto
FROM pagos p
JOIN autores a ON p.id_autor = a.id_autor
ORDER BY p.fecha_pago;
\end{lstlisting}

\captura{consulta_10.1.png}{Resultados de la consulta 10, sentencia 1 en MySQL Workbench.}{fig:c10a}

\clearpage
\subsection*{Sentencia 2: total pagado a cada autor}

\begin{lstlisting}
SELECT a.nombre, SUM(p.monto) AS total_pagado
FROM autores a
JOIN pagos p ON a.id_autor = p.id_autor
GROUP BY a.id_autor, a.nombre
ORDER BY total_pagado DESC;
\end{lstlisting}

\captura{consulta_10.2.png}{Resultado de la consulta 10, sentencia 2 en MySQL Workbench.}{fig:c10b}

\end{document}
